\documentclass[10pt,a4paper]{article}
\usepackage[utf8]{inputenc}
\usepackage{amsmath}
\usepackage{amsfonts}
\usepackage{comment}
\usepackage{amssymb}
\usepackage{graphicx}
\usepackage[left=2cm,right=2cm,top=2cm,bottom=2cm]{geometry}
\usepackage{listings}
\usepackage{color}
\definecolor{munsell}{rgb}{0.0, 0.5, 0.69}
\definecolor{dkgreen}{rgb}{0,0.6,0}
\definecolor{gray}{rgb}{0.5,0.5,0.5}
\definecolor{mauve}{rgb}{0.58,0,0.82}
\definecolor{minas}{RGB}{0.244, 0.172, 0.36}
\definecolor{PUJ}{RGB}{44, 86, 151}
\definecolor{PUJ2}{RGB}{43, 93, 156}
\definecolor{PUJ3}{RGB}{20, 52, 107}
\definecolor{cyandk}{rgb}{0.0, 0.72, 0.92}

\author{Iván Andrés Trujillo }

%% code information
\lstset{frame=tb,
  language=Python,
  aboveskip=3mm,
  belowskip=3mm,
  showstringspaces=false,
  columns=flexible,
  basicstyle={\small\ttfamily},
  numbers=none,
  classoffset=1,
  morekeywords={True,False}, keywordstyle=\color{munsell}, 
  classoffset=0, 
  keywordstyle=\color{blue},  
  commentstyle=\color{dkgreen},
  stringstyle=\color{PUJ3},
  numberstyle=\tiny\color{gray},
  breaklines=true,
  breakatwhitespace=true,
  tabsize=4,
}

\begin{document}

\newcommand{\HRule}{\rule{\linewidth}{0.5mm}}

%----------------------------------------------------------------------------------------
%	LOGO SECTION
%----------------------------------------------------------------------------------------

\includegraphics[width = 4cm]{pujshield.eps}\\[0.5cm] 

\begin{center} 
\textbf{\LARGE Introduction to regular expressions}\\[0.2cm]
\emph{\LARGE Notes of class}\\[0.3cm] 
\emph{Iván Andrés Trujillo Abella} \\
\textsc{\Large 
}\\[0.2cm] 
\textsc{\large Facultad de ingeniería}\\[0.5cm] 
\HRule \\[0.4cm]
\end{center}
\vspace{1cm}


\section{Freddy investigation}
ASCII ?
UNICODE?




\section{Regular expressions}
Find patterns in a string ( search pattern).

are related with wildcards



\subsection{Quantifiers}

How many times can occur 




\begin{verbatim}
\b[A-Z0-9._%+-]+@[A-Z0-9.-]+\.[A-Z]{2,}\b
\end{verbatim}


\begin{comment}
thereof: del mismo.
sparingly:escasamente
collides: 
comes down:
stumbled:
\end{comment}


\subsection{special meaning}

Metacharacteres
\begin{itemize}
\item \textbackslash
\item .
\item \textasciicircum
\item \textbar
\item ?
\item *
\item +
\item ()
\item $[]$
\end{itemize}

override the special meaning with \textbackslash. 

important thing:
windows end files with \textbackslash r\textbackslash n  and unix \textbackslash n.


$[]$ match only a single iput $a,b,..$ for instance:
\begin{verbatim}
ca[tr]^
\end{verbatim}
could math car or cat..


whith - we list a sequence of characteres.

\begin{verbatim}
[A-Z]
\end{verbatim}


we can uses  \textasciicircum to avoid inside the square bracket the character, for instance:
\begin{verbatim}
le[^o]n
\end{verbatim}

\begin{verbatim}
leon
len
\end{verbatim}

\subsection{.}
Match any character except a new line \textbackslash n.


\subsection{position}{ \textasciicircum , \$}


\textasciicircum  at the start the string

\$ to the end of string

\subsubsection{word boundary}

\textbackslash b is a combination of  and  \textasciicircum\$. boundary word.

we can uses the boundary to find full words, what happen when the word finis with (!,?, - or whit space)?.
\begin{verbatim}
regex = \bwords\b
words compound language
sentences are composesd of words_
words?
words"
\end{verbatim}

\textbacklash B find when the patter no it is and start or end the word.





\subsection{consider}
\begin{verbatim}
\w  = [A-Za-z0-9_]
\d = [0-9]
\D = [^0-9]
\s = any space, tab or newline.
\end{verbatim}



\subsection{or}

\begin{verbatim}
regex = (ML|deep learning) techniques
ML deep learning techniques
ML techniques
deep learning techniques 
ML is a technique
\end{verbatim}


\begin{verbatim}
regex = ML|deep learning techniques
in a text could be ML techniques 
or deep learning techniques becuase 
\end{verbatim}


\section{repetition}
? indicating optional, match zero or more.
\begin{verbatim}
regex = "?optional
optional\n)
"optional
\end{verbatim}


* zero or more times
+ once or more
\{ int \} the number of repetitions allowed.
match only number of four digits:
\begin{verbatim}
regex = \b[0-9]{4}\b
1234
4566
89091
\end{verbatim}


\begin{verbatim}
regex = \b[1-9][0-9]{int,int}\b
\end{verbatim}
integers define the range of limitations.

\begin{verbatim}
regex = \b[a-z]{2,5}\b
last 
focus on simple 
\end{verbatim}

We can say that the number of repetitions are in the last code we can say that \{ at least , at most \}.
if 'at most' not is defined then assume \{at least , \} that the uppermost boundary is indefinite. 



\section{take in mind}
*,+,? are greedy, they match all text.

inside the square brackets not is neccesary escape quantifiers.

? symbol could be used to the end of a quantifier to make a not greedy search.
also could be used with \{ \}, in this case match only few as possible.

\begin{verbatim}
regex = carto{1,n}?
cartooon
\end{verbatim}

in a general case only match \{ at least 


\begin{lstlisting}
regex = 'c.*n'
regex = 'c.*?n
\end{lstlisting}



\section{()}
start and the end of a group



\section{Another examples}
\begin{verbatim}
regex = cart(o){1,}n
cartoon 
cartoooooon 
\end{verbatim}

\section{\textbar}
not is greedy.


\section{$[]$}

the square brackets is a set.


\begin{lstlisting}
'[A-Z][a-z]+', 'one uppercase followed by lowercase'
\end{lstlisting}


\section{Excercises}
match the vowels.












\section{search and findall methods in python}
\begin{lstlisting}
import re
re.match(regex, string)
\end{lstlisting}


\begin{lstlisting}

\end{lstlisting}






\begin{lstlisting}
re.search(regex, string) # the first location
re.findall(regex, string) #return an array 
\end{lstlisting}



\begin{lstlisting}
string = " Here it is important recognize some emails for instance IA@gmail.com  or maybe support@somemail.es"
re.findall('[a-z]+@[a-z]+\.[a-z]*',texto)
\end{lstlisting}



\subsection{avoid}
The symbol caret allow us not search a set of characters.
\begin{verbatim}
[^abc]
\end{verbatim}



\begin{itemize}
\item confidence interval
\item difference paired data
\item test population proportion
\item comparing two independent means
https://www.coursera.org/learn/inferential-statistical-analysis-python#syllabus

\end{itemize}




\begin{re.sub()}
replace letters by empty.
\begin{lstlisting}
re.sub(regex, to_replace, count)
#count how many replacements will be 
\end{lstlisting}

Could be useful the following code:
\begin{lstlisting}
import re

def repl_func(match):
  if match == True:
    return " "

string = "Hello! How are you?! Where have you been?!"
new_string = re.sub(r"""[!?'".<>(){}@%&*/[/]""", repl_func, string)
print(new_string)
\end{lstlisting}




\section{Fuzzy string matching}

if $a,b$ are string the edit distance is the number of operations of edition to transform $a$ in $b$.

We could think in the following operations over \xi (empty string and y )
\begin{itemize}
\item  insertion
\item  deletion
\item  substitution
\end{itemize}




\subsection{Edit distance}
\subsection{Background}

Similarity of DNA sequences:


\begin{comment}
Los escapes son indicados por el prefijo del caracter con una barra invertida (\). Desafortunadamente, una barra invertida debe ser escapada en sí misma en cadenas de Python, y eso resulta en expresiones difíciles de leer. Usar cadenas raw, que se crean al ponerle un prefijo al literal valor con r, elimina este problema y mantiene la legibilidad.


https://rico-schmidt.name/pymotw-3/re/index.html
good documentation

\end{comment}


\end{document}
